% ============================================================
% main.tex - Day 4 packet, "The Pilot Inquiry".
% ChemicalQDevice daily funding action packet, repository v4.9.0.
% Deposited in funding/auto-fund-2/01Oct26/packet.
%
% WHAT THIS DOCUMENT IS.  The fourth packet of the block.  It answers an inquiry from
% the National Science Foundation about applying to a multi-million-dollar pilot
% program, states the pilot's published eligibility gate and the company's place
% against it, and sets out the Project Pitch that opens the only route to it.
%
% STRUCTURE.  Seven sections, one .tex file each, per Rule 6.  Three figures,
% five tables.  Every figure is TikZ; no raster is read or written.
%
% SPACING INVARIANT.  \vspace{-0.60cm} before every caption, giving 7.44pt from
% the last rule to the first caption line, identically for every float.
%
% COMPILE ON OVERLEAF WITH pdfLaTeX:
%   pdflatex main -> bibtex main -> pdflatex main -> pdflatex main
% ============================================================
\documentclass[11pt]{article}
\usepackage{fundstyle}

\title{The Pilot Inquiry: The Published Gate to a Multi-Million-Dollar NSF Pilot, and the Project Pitch That Opens the Route}
\author{Kevin Kawchak}
\date{October 1, 2026}

\fundhdr{The Pilot Inquiry}{Day 4 of 5}

\begin{document}
\thispagestyle{plain}

\fundmast{ChemicalQDevice \;|\; Daily Funding Action Packet \;|\; Day 4 of 5}%
{The Pilot Inquiry}%
{The National Science Foundation asked whether the company would apply to a
multi-million-dollar pilot program. This packet answers exactly: the pilot's published
gate, where the company stands against it, and the Pitch that takes the first step today.}

\vspace{0.15em}

\begin{center}
\appbadge{Day 4 of 5} \enspace \appbadgel{Repository v4.9.0} \enspace
\appbadgel{3 figures} \enspace \appbadgel{5 tables} \enspace
\appbadgel{Pitch first, pilot later}
\end{center}

\vspace{-0.15em}

\funddeck{The pilot}{\$20M}{Six to eight awards of \$1M to \$3.75M}%
{The gate}{Phase II}{An NSF award the company does not yet hold}%
{The first step}{\$305K}{Phase I, opened by the Pitch submitted today}

\vspace{-0.1em}
\apprule{0.9pt}

\begin{center}\small
CEO Kevin Kawchak\,\orcidicon{0009-0007-5457-8667}, ChemicalQDevice \;|\;
\href{mailto:kevinkawchak@gmail.com}{kevinkawchak@gmail.com}\\
\href{https://orcid.org/0009-0007-5457-8667}{0009-0007-5457-8667} \;|\;
San Diego, California
\end{center}

\vspace{-0.25em}
\apprule{0.5pt}

\noindent\footnotesize\textit{Independent research paper and practical adoption
guide. It is not medical or regulatory advice and is not endorsed by the FDA,
NIH, HHS, an IRB, ICH, or any sponsor. All figures derive from the author's
repository sources and are illustrative unless tied to a cited
reference.}\normalsize\par

\vspace{0.12em}

\noindent\footnotesize\textbf{Disclaimer:} \textit{This work is independent and
is not endorsed or sponsored by any trial sponsor, CRO, site, IRB, regulator, or
medical society; and was adapted using Claude Code Opus 5.5.}\normalsize\par

\vspace{0.12em}

{\sloppy\noindent\footnotesize Repository (v4.9.0) is at
\url{https://github.com/kevinkawchak/physical-ai-oncology-trials}.
The paper files are deposited in
\href{https://github.com/kevinkawchak/physical-ai-oncology-trials/tree/main/funding/auto-fund-2}{/funding/auto-fund-2}.
No digital object identifier is asserted for this document.\normalsize\par}

\renewcommand{\keywordname}{\textbf{Keywords}}
\keywords{National Science Foundation; SBIR Phase I; Project Pitch; Commercialization Readiness Pilot; Strategic Breakthrough; Language model verification; Robotic surgery; Pancreatic ductal adenocarcinoma}

\vspace{0.25em}
\apprule{0.6pt}

\setcounter{tocdepth}{1}
\renewcommand{\contentsname}{\color{accentcol}Table of Contents}
{\setlength{\parskip}{0pt}\footnotesize\linespread{0.94}\selectfont\tableofcontents}

\vspace{-0.1em}
\apprule{0.6pt}

% The body follows the contents on the cover page rather than after a page
% break, so the cover carries no block of empty space below its contents.
% ---------- Body: one \input per section (Rule 6) ----------
% sec-00: front matter.  Abstract and the reading path.
\section*{Abstract}
\addcontentsline{toc}{section}{Abstract}

The National Science Foundation has asked whether ChemicalQDevice would apply to
a multi-million-dollar pilot program. In September 2026 the agency announced a
\$20 million, two-year Commercialization Readiness Pilot, run by the National
Commercialization and Translation Institute at the University of Central
Florida, that will select six to eight companies for milestone-based awards of
\$1 million to \$3.75 million each \cite{nsfncti2026,nctiucf2026}. Its competition
is open to active or recent NSF SBIR and STTR Phase II and IIB awardees. The
agency's Strategic Breakthrough tier, of up to \$30 million with a one-to-one
match, is likewise open only to Phase II awardees \cite{nsf26510}. The company
holds no NSF award, and so meets neither gate today.

This packet answers the inquiry exactly. It thanks the agency, states the gate
and the company's place against it, asks whether the inquiry contemplates a
route the company has missed, and submits the Project Pitch that is the first
step on the published route: Pitch, Phase I, Phase II, and only then the pilot.
NSF's SBIR program funds technical innovation and does not fund clinical
trials, so the Pitch proposes the part of the program NSF can fund: a governed,
advisory-only language model layer for robotic oncologic surgery, with
verification by independent models and a public audit trail. The robotic
Phase 1 itself stays with the National Institutes of Health and private capital,
and the packet shows how the program is split so that no dollar is asked for
twice.

\subsection*{How to read this}

A reader with five minutes takes \S1, the inquiry, the gate and the funding
ladder, and \S2, the day's actions with the path from Pitch to proposal.

A reviewer asking what NSF is being asked to fund takes \S3, the split between
NSF, NIH and private capital, and \S4, the four Phase I objectives, each with a
measure and a threshold.

A funder takes \S5, the merit review criteria against the program's evidence,
including the quantitative record behind the robotic Phase 1. \S6 carries the
positioning constraints, the method, and the references, every one of which is
a clickable link.

% sec-01: the inquiry, the gate, and the funding ladder.  Table 16, Figure 10
% (d2-type layered stack).
\section{The Inquiry, and the Gate in Front of the Pilot}

\subsection{What was asked, and what the answer must be}

An inquiry to apply is an honor, and it is also a question with an exact
answer. The pilot the inquiry concerns is new: NSF announced it in September
2026 as a way to carry deep technologies from its small business portfolio
toward market readiness, with the Institute at the University of Central Florida
selecting the companies and managing their milestones \cite{nsfncti2026}. The
gate to it is published, and the company does not yet pass it. The answer
therefore does three things and no more: it thanks the agency, it states the
gate and the company's place against it in writing, and it asks whether the
inquiry had a different route in mind. The letter to the pilot's principal
investigator asks the same question of the Institute \cite{nctiucf2026}.

\subsection{The tiers and pilots, side by side}

Table 16 sets every NSF tier and pilot the inquiry could concern beside its
amount, its eligibility, and the company's position today. The amounts are
those published in the current solicitations and the pilot's announcement
\cite{nsf26510,nsf26511,nsfncti2026}. Only one row is open to the company now,
and it carries no money: the Project Pitch, which NSF requires before any
Phase I proposal.

\begin{apptable}
\begin{tabularx}{\textwidth}{>{\raggedright\arraybackslash}p{3.9cm} >{\raggedright\arraybackslash}p{4.2cm} >{\raggedright\arraybackslash}p{3.8cm} >{\raggedright\arraybackslash}X}
\headrow \thead{Tier or pilot} & \thead{Amount} & \thead{Who is eligible} & \thead{The company}\\
\midrule
Project Pitch & None; four short fields & Any small business & Submitted today \\
SBIR Phase I & \$305K; 6 to 18 months & Invited after a Pitch & Once invited \\
SBIR Phase II & \$1.25M; about 24 months & Phase I awardees & After Phase I \\
Phase IIB supplement & \$50K to \$500K, with a match & Phase II awardees & Ledger at \$0 \\
NCTI Readiness Pilot & \$1M to \$3.75M each, of \$20M & Recent Phase II or IIB & Not yet eligible \\
Strategic Breakthrough & Up to \$30M; 1 to 1 match & Phase II awardees & Not yet eligible \\
Instrumentation, 26-511 & Share of a \$40M emphasis & Instrument proposers & Outside scope \\
\end{tabularx}
\end{apptable}
\vspace{-0.60cm}
\tabcap{Table 16. The NSF tiers and pilots the inquiry could concern, each with its amount, who\\
is eligible, and the company's place today; only the Project Pitch is open to it now.}

\subsection{The same tiers, drawn as a ladder}

Figure 10 draws Table 16 as the stack it really is. Each layer is reached only
through the one beneath it, and the gate between two layers is an award, a
match, or an invitation. The company stands on the lowest layer. The pilot the
inquiry concerns sits two gates above the first award the company could win.

\begin{appfloat}[!tb]
\begin{appfig}[d2-type / layered stack, each layer gated by the one below]
% Six layers bottom to top, each a D2 container 11.0 cm wide with the tier on
% the left and its amount on the right.  The gate that opens a layer is set in
% the right-hand column at the boundary beneath it.  The company marker is the
% one accent layer, at the bottom.
\tikzset{d2tier/.style={text width=48mm,minimum height=7.2mm,align=left,anchor=west}}
\node[d2key,d2tier,font=\scriptsize\sffamily\bfseries] at (0.00,0.00) {Project Pitch};
\node[d2key,d2tier] at (5.30,0.00) {No funding; four fields; two a year};
\node[d2soft,d2tier,font=\scriptsize\sffamily\bfseries] at (0.00,1.05) {SBIR Phase I};
\node[d2soft,d2tier] at (5.30,1.05) {Up to \$305K, 6 to 18 months};
\node[d2soft,d2tier,font=\scriptsize\sffamily\bfseries] at (0.00,2.10) {SBIR Phase II};
\node[d2soft,d2tier] at (5.30,2.10) {Up to \$1.25M, about 24 months};
\node[d2gray,d2tier,font=\scriptsize\sffamily\bfseries] at (0.00,3.15) {Phase IIB supplement};
\node[d2gray,d2tier] at (5.30,3.15) {\$50K to \$500K, with a match};
\node[d2gray,d2tier,font=\scriptsize\sffamily\bfseries] at (0.00,4.20) {Commercialization Readiness Pilot};
\node[d2gray,d2tier] at (5.30,4.20) {\$1M to \$3.75M each; 6 to 8 firms};
\node[d2gray,d2tier,font=\scriptsize\sffamily\bfseries] at (0.00,5.25) {Strategic Breakthrough};
\node[d2gray,d2tier] at (5.30,5.25) {Up to \$30M, 1 to 1 match};
\begin{scope}[on background layer]
\foreach \y in {0.00,1.05,2.10,3.15,4.20,5.25}{
  \draw[fundmid,line width=0.5pt,rounded corners=3pt,fill=fundwhite] (-0.15,\y-0.47) rectangle (10.75,\y+0.47);}
\end{scope}
\foreach \y/\g in {0.52/invitation after the Pitch, 1.57/a Phase I award, 2.62/a Phase II award, 3.67/Phase II or IIB, 4.72/Phase II and a match}{
  \node[font=\tiny\sffamily,text=fundgrayd,anchor=west] at (11.25,\y) {gate: \g};
  \draw[d2edged] (11.15,\y-0.30) -- (11.15,\y+0.30);}
\node[d2key,text width=24mm,minimum height=7.2mm] (here) at (13.35,0.00) {The company is here};
\draw[d2edgeb] (here.west) -- (10.80,0.00);
\legkey{0.00}{-0.95}{fundaccent}{Open to the company today}
\legkey{4.30}{-0.95}{fundpale}{Reached by winning the layer below}
\legkey{9.05}{-0.95}{fundgrayl}{Needs a Phase II award first}
\ptitle{-0.10}{6.20}{Each layer opens only from the one beneath it}
\end{appfig}
\vspace{-0.60cm}
\figcaption{Figure 10. The NSF tiers and pilots as a layered stack, each opened only by an award\\
or a match on the layer below, with the company on the Pitch layer at the bottom.}
\end{appfloat}

\subsection{The standing fact}

The trial behind the program is built around daraxonrasib, which the FDA
approved on August 26, 2026 as Rasonque for metastatic pancreatic adenocarcinoma
after prior systemic therapy \cite{fdarasonque2026,fdadisco2026,revmedapproval2026}.
The approval followed Revolution Medicines' Phase 1/2 program and its Phase 3
RASolute 302 trial \cite{ctgov001,ctgov302,rasolute302}. ChemicalQDevice's AI
papers on the molecule were developed independently of that program and
simultaneously with it, from the June 2025 identification study
\cite{daraxident} to the Phase 1 protocol and investigational new drug
application of 2026 \cite{onpremwhippl,phase1ind}. The perioperative use the
company's protocol proposes remains investigational, and a separate adjuvant
Phase 3 trial of the molecule has been enrolling since December 2025
\cite{ctgov304}.

% sec-02: the action register.  Table 17, Figure 11 (plantuml-type activity
% diagram with a fork, a join, and a connector).
\section{The Action Register}

Table 17 lists the day's seven actions in the order they are done. The answer to
the inquiry goes first, because every other letter refers to it. The Pitch waits
on the organization record, and the letter to NIH waits on the Pitch, because it
reports the Pitch as submitted.

\begin{apptable}
\begin{tabularx}{\textwidth}{>{\raggedright\arraybackslash}p{0.5cm} >{\raggedright\arraybackslash}p{6.4cm} >{\raggedright\arraybackslash}p{2.6cm} >{\raggedright\arraybackslash}X}
\headrow \thead{No} & \thead{Action} & \thead{File} & \thead{Waits on}\\
\midrule
1 & Answer the NSF inquiry & {\footnotesize\ttfamily email-01} & Nothing \\
2 & Ask the pilot's gate of its investigator & {\footnotesize\ttfamily email-02} & Row 1 sent \\
3 & Confirm the Research.gov record & {\footnotesize\ttfamily form-02} & An active SAM.gov record \\
4 & Submit the Project Pitch & {\footnotesize\ttfamily form-01} & Row 3 \\
5 & Ask about NAIRR Pilot resources & {\footnotesize\ttfamily email-03} & Nothing \\
6 & Disclose the Pitch to NIH & {\footnotesize\ttfamily email-04} & Row 4 submitted \\
7 & Open the match ledger & {\footnotesize\ttfamily capital-04} & Nothing \\
\end{tabularx}
\end{apptable}
\vspace{-0.60cm}
\tabcap{Table 17. The day's seven actions in the order they are done, each with the file that\\
carries it and what it waits on; the letter to NIH follows the Pitch it discloses.}

\subsection{From the Pitch to a proposal}

Figure 11 draws the route from the answer to a proposal as an activity diagram.
Two branches run in parallel and must join before the Pitch is submitted: the
registration branch, which confirms the company's federal entity record and its
Research.gov organization \cite{researchgov2026}, and the drafting branch, which
writes the four Pitch fields and measures each against its character limit. The
diagram breaks at a connector and resumes at the right. NSF then invites a full
proposal or declines to \cite{nsf26510}. A declined Pitch is revised and
submitted again, and the solicitation allows two Pitches a year.

\begin{appfloat}[!tb]
\begin{appfig}[plantuml-type / activity diagram with a fork, a join, and a connector]
% Left block: the answer, a fork into two branches, the join, and the Pitch,
% ending at connector A.  Right block: connector A, the NIH disclosure, the
% one decision, and the stop, with a note on what an award would open.
\node[umlinit] (st) at (3.50,0.30) {};
\node[umlbox,text width=46mm] (a1) at (3.50,-0.35) {Answer the NSF inquiry, letter 1};
\node[umlbar,minimum width=70mm] (fk) at (3.50,-1.00) {};
\node[umlbox,text width=40mm] (b1) at (1.25,-1.65) {Confirm SAM.gov and the UEI};
\node[umlbox,text width=40mm] (b2) at (1.25,-2.60) {Research.gov organization};
\node[umlbox,text width=40mm] (b3) at (1.25,-3.55) {Assign the PI and AOR roles};
\node[umlbox,text width=40mm] (c1) at (5.75,-1.65) {Draft the four Pitch fields};
\node[umlbox,text width=40mm] (c2) at (5.75,-2.60) {Check each field's limit};
\node[umlbox,text width=40mm] (c3) at (5.75,-3.55) {Confirm no field proposes a trial};
\node[umlbar,minimum width=70mm] (jn) at (3.50,-4.20) {};
\node[umlkey,text width=46mm] (sb) at (3.50,-4.85) {Submit the Project Pitch};
\node[circle,draw=fundink,line width=0.6pt,fill=fundwhite,minimum size=4.2mm,
  inner sep=0pt,font=\tiny\sffamily\bfseries] (ca) at (3.50,-5.55) {A};
\draw[umlarrow] (st) -- (a1);
\draw[umlarrow] (a1) -- (fk);
\draw[umlarrow] (fk.south -| b1) -- (b1);
\draw[umlarrow] (fk.south -| c1) -- (c1);
\draw[umlarrow] (b1) -- (b2);
\draw[umlarrow] (b2) -- (b3);
\draw[umlarrow] (c1) -- (c2);
\draw[umlarrow] (c2) -- (c3);
\draw[umlarrow] (b3) -- (jn.north -| b3);
\draw[umlarrow] (c3) -- (jn.north -| c3);
\draw[umlarrow] (jn) -- (sb);
\draw[umlarrow] (sb) -- (ca);
\node[circle,draw=fundink,line width=0.6pt,fill=fundwhite,minimum size=4.2mm,
  inner sep=0pt,font=\tiny\sffamily\bfseries] (cb) at (10.60,0.30) {A};
\node[umlbox,text width=46mm] (d1) at (10.60,-0.35) {Disclose the Pitch to NIH, letter 4};
\node[diamond,aspect=2.4,draw=fundink,line width=0.6pt,fill=fundwhite,inner sep=1pt,
  text width=26mm,font=\tiny\sffamily,align=center] (dc) at (10.60,-1.55) {Invited to submit a proposal?};
\node[umlbox,text width=40mm] (p1) at (10.60,-2.95) {Proposal by the next deadline};
\node[umlbox,text width=40mm] (p2) at (10.60,-3.90) {Merit review on NSF's criteria};
\node[umlbox,text width=24mm] (r1) at (14.35,-2.95) {Revise and pitch again};
\node[diamond,draw=fundink,line width=0.6pt,fill=fundwhite,minimum size=3.6mm,inner sep=0pt] (mg) at (10.60,-4.70) {};
\node[umlfinal] (en) at (10.60,-5.45) {};
\fill[fundink] (en) circle (0.9mm);
\node[umlnote,text width=27mm] (nt) at (14.10,-5.45) {Only a Phase II award opens the pilot and Strategic Breakthrough};
\draw[umlarrow] (cb) -- (d1);
\draw[umlarrow] (d1) -- (dc);
\draw[umlarrow] (dc.south) -- node[umlguard,right,pos=0.35] {[yes]} (p1.north);
\draw[umlarrow] (dc.east) -| node[umlguard,above,pos=0.25] {[no]} (r1.north);
\draw[umlarrow] (p1) -- (p2);
\draw[umlarrow] (p2) -- (mg);
\draw[umlarrow] (r1.south) |- (mg.east);
\draw[umlarrow] (mg) -- (en);
\draw[umldash,-] (nt.west) -- (en.east);
\ptitle{-0.75}{0.95}{Two branches join before the Pitch}
\end{appfig}
\vspace{-0.60cm}
\figcaption{Figure 11. The route from the answer to a proposal, with registration and drafting\\
run in parallel and joined before the Pitch, then one decision that NSF alone makes.}
\end{appfloat}

\subsection{What each letter asks}

Letter 1 thanks NSF's SBIR program for the inquiry, states the gate, and asks
whether the inquiry contemplates a route the company has missed. It copies the
Policy Office, which owns the rules the gate is written in, and the Research.gov
help desk, which would resolve any registration fault \cite{nsfseedcontact,nsfpappg}.
Letter 2 asks the pilot's principal investigator at the University of Central
Florida to confirm the gate in writing, and asks for nothing else
\cite{nctiucf2026}. Letter 3 asks the NAIRR Pilot whether a small business can
request shared resources for verification research on synthetic data only
\cite{nsfnairr2026,nairrpilot2026}. Letter 4 tells the NIH SBIR program office
and the National Cancer Institute's SBIR team that an NSF Pitch has been
submitted on the verification work, so that any later NIH application discloses
it and neither agency pays for the same work twice
\cite{nihsbirseed,ncisbir2026,sbapolicydirective}.

\subsection{A note on the date}

October 1 is the first day of the federal fiscal year. If appropriations have
lapsed, NSF staff may not answer for a time, and a portal may accept a
submission it cannot yet process. The Pitch can still be submitted. Nothing else
is followed up until three business days after the agency reopens.

\actionbox{The one approval this day asks for}{Approve answering the NSF inquiry,
testing the pilot's published eligibility gate in writing with its principal
investigator, and submitting the Project Pitch that opens the only route to it.
Approve disclosing the Pitch to NIH on the day it is submitted.}

% sec-03: how the program is split between NSF, NIH and private capital.
% Table 18.
\section{How the Program Is Split}

\subsection{The rule that shapes the split}

NSF's SBIR program funds technical innovation with commercial potential, and it
does not fund clinical trials \cite{nsf26510}. The National Institutes of Health
fund clinical research, including small business work that leads to a trial
\cite{nihsbirseed,ncisbir2026}. The program's private capital sits behind a
firewall that gives no investor any say in trial conduct, and pays for what
neither agency pays for \cite{capitalplan2026}. The split in Table 18 follows
those three facts and nothing else. No component is marked yes in two public
columns.

\begin{apptable}
\begin{tabularx}{\textwidth}{>{\raggedright\arraybackslash}p{5.3cm} >{\raggedright\arraybackslash}p{3.2cm} >{\raggedright\arraybackslash}p{3.5cm} >{\raggedright\arraybackslash}X}
\headrow \thead{Component} & \thead{NSF SBIR} & \thead{NIH SBIR} & \thead{Private capital}\\
\midrule
Verification method for the model & Yes, the Pitch & No & No \\
Boundary tests and audit trail & Yes & Cites the NSF result & No \\
Timing tests, hardware in the loop & Yes & No & No \\
Protocol, IND, site activation & No & Yes & In part \\
Drug, surgery, and follow-up & No & Yes, at the site & In part \\
Robot platform, site build-out & No & In part & Yes \\
Commercializing the method & After Phase II & No & Yes \\
\end{tabularx}
\end{apptable}
\vspace{-0.60cm}
\tabcap{Table 18. The program's seven components split between NSF, NIH and private capital,\\
with no component funded twice; the verification work is the part proposed to NSF.}

\subsection{Why the scopes cannot overlap}

The NSF work ends where a patient begins. It produces a verification method, a
test suite and measured failure rates on synthetic cases generated from
published parameters. The NIH work begins where a patient begins: it uses the
verified layer inside a trial that belongs to a site investigator, under an
investigational new drug application filed under 21 CFR Part 312
\cite{ecfr2026part312}. If both were funded, the NIH work would cite the NSF
results rather than repeat them, and each application would disclose the other,
as the SBIR policy directive requires \cite{sbapolicydirective}. Letter 4 tells
NIH so in writing, with NSF copied.

\subsection{The money on each side}

An NSF Phase I is worth up to \$305,000 \cite{nsf26510}. The company's plan
carries an NIH Phase I of \$306,000 and an NIH Phase II of \$1,300,000, for an
SBIR route of \$1,606,000, and places \$5,900,000 of private capital behind the
firewall against the \$2,104,000 delta the route does not buy
\cite{capitalplan2026}. The NSF and NIH Phase I figures are nearly equal and pay
for different work, and neither is counted twice anywhere in the company's
documents. The private figure is a plan, not a completed raise, and the day's
match ledger records its opening balance as zero.

\subsection{What the split means for the trial}

None of the robotic Phase 1 is proposed to NSF. What NSF would fund is the
evidence that the advisory layer inside that trial stays inside its boundaries.
That is the evidence an institutional review board and the FDA will ask for
before the layer is allowed near a patient \cite{fdaaiguidance2025}, and it is
the evidence a clinical funder will ask for before paying for the trial it sits
inside. The NSF work therefore does not compete with the NIH route. It is the
step that makes the NIH route fundable.

% sec-04: the Phase I technical objectives.  Table 19.
\section{The Phase I Technical Objectives}

\subsection{What the Pitch proposes}

The Pitch proposes a governed, advisory-only language model layer for robotic
oncologic surgery, and a verification method that makes its behavior checkable.
The layer runs on premises at a clinical site. It produces text recommendations
to the operating surgeon and nothing else: it holds no write credential to the
data capture system, no route to the robot control network, and no ability to
act on equipment \cite{onpremwhippl,phase1ind}. Every output is checked by two
independent models from different developers before it is shown, a practice the
company already follows in its own work \cite{tripartisan2026}, and every check
is written to an audit trail that can be replayed.

The innovation is the combination of three things that clinical AI systems
usually treat separately: a hard architectural boundary, enforced by network
segmentation and credential design rather than by instructions to the model;
cross-model verification, in which disagreement between independent models is a
signal to withhold advice rather than noise; and a credibility assessment framed
on the ASME V\&V 40 approach to computational models \cite{asmevv40}, extended
from simulation to advisory text.

\subsection{Five measures, each with a threshold}

Table 19 states the Phase I as five measures. The first four are the technical
risks the Pitch names; the fifth restates the physical limits of the deposited
protocol that the layer must never touch. Each row has a threshold a reviewer
can check at the end of the award.

\begin{apptable}
\begin{tabularx}{\textwidth}{>{\raggedright\arraybackslash}p{2.9cm} >{\raggedright\arraybackslash}p{5.0cm} >{\raggedright\arraybackslash}p{4.6cm} >{\raggedright\arraybackslash}X}
\headrow \thead{Objective} & \thead{Measure} & \thead{Threshold} & \thead{Source}\\
\midrule
Boundary integrity & Open paths across the boundary, tested on every build & Zero write credentials; no route to robot control & Protocol, IND \\
Verification yield & How often verifiers disagree, and flag a seeded error & A withholding threshold set from measured rates & Company method \\
Credibility & An ASME V\&V 40 battery, extended to advisory text & Twin at 81.9 over 55 tests; extension with limits & Deposited twin \\
Timing & Latency against each safety stop, on a hardware bench & No effect on 3 ms arm or 500 ms system stops & Protocol \\
Physical limits & Any advisory path to a force or vessel control & Zero paths to the 3 N, 18 N and five-vessel limits & Protocol \\
\end{tabularx}
\end{apptable}
\vspace{-0.60cm}
\tabcap{Table 19. The Phase I as five measures, each with a threshold a reviewer can check at\\
the end of the award and the deposited source it comes from; no patient data is used.}

\subsection{What the Phase I would produce}

A Phase I would produce three things: a verification method, the test suite that
exercises it on every build, and a public report of measured rates with their
limitations. The main technical risks are building synthetic cases realistic
enough to expose failure, choosing verifier models that fail independently
rather than together, and measuring rare events without patient data. Each is a
question of engineering and measurement, which is the kind of question NSF
funds.

\subsection{What is not proposed}

No clinical trial is proposed, and no patient data is used, identifiable or
de-identified. The deployed clinical model stays on premises at the site, and
any shared computing resource would see only synthetic cases and the company's
own code. The robotic Phase 1 the layer would one day sit inside remains an
open-label, single-arm, 3+3 dose escalation of perioperative daraxonrasib at
160, 220 and 300 mg, with a 28-day dose-limiting toxicity window, in up to 18
treated participants undergoing a staged eight-arm robotic
pancreaticoduodenectomy with 56 degrees of freedom and 640 sensor channels
\cite{onpremwhippl}. That trial belongs to a site investigator and to NIH and
private funding, as \S3 sets out.

% sec-05: merit review, where the work would run, and the quantitative record.
% Table 20, Figure 12 (diagrams (python)-type clustered infrastructure).
\section{The Merit Review Map}

NSF reviews every proposal against two criteria, Intellectual Merit and Broader
Impacts, and judges SBIR proposals on their commercial impact as well
\cite{nsfpappg,nsf26510}. Table 20 maps each question a reviewer will ask to the
evidence the company can put in front of that reviewer today, and names the
evidence it does not yet have. The last column is as important as the third: a
reviewer trusts a map that shows its own gaps.

\begin{apptable}
\begin{tabularx}{\textwidth}{>{\raggedright\arraybackslash}p{2.5cm} >{\raggedright\arraybackslash}p{3.6cm} >{\raggedright\arraybackslash}p{5.4cm} >{\raggedright\arraybackslash}X}
\headrow \thead{Criterion} & \thead{What NSF asks} & \thead{The evidence today} & \thead{Not yet in hand}\\
\midrule
Intellectual Merit & Advances knowledge? & Model disagreement used as a signal to withhold & Measured rates \\
Intellectual Merit & Well reasoned? & A hard boundary, not instructions; ASME V\&V 40 & An outside review \\
Intellectual Merit & Team qualified? & A dated public record since June 2025 & Staff, on award \\
Intellectual Merit & Resources adequate? & On-premises compute; a NAIRR request & A NAIRR reply \\
Broader Impacts & Benefit to society & Models used only once shown to stay in bounds & Phase I evidence \\
Broader Impacts & Workforce & Eleven planned roles, six needing no license & A first hire \\
Commercial impact & Who pays, and why? & Centers, sponsors and robot makers, for a checkable boundary & A paying customer \\
\end{tabularx}
\end{apptable}
\vspace{-0.60cm}
\tabcap{Table 20. The seven questions an NSF reviewer asks under the three criteria, with the\\
evidence the company can show today and the evidence it does not yet have in hand.}

\subsection{Where compute, data and review would sit}

The resources question is answered by a drawing rather than a sentence. Figure 12
places every component in one of three clusters: the clinical site, where the
model runs and where patient data would stay; the shared NAIRR Pilot resources,
which would see only synthetic cases and the company's code if a request were
granted \cite{nairrpilot2026}; and the public record, which receives the method
and the measured rates. The robot control network is a cluster of its own inside
the site, and the model has no route into it.

\begin{appfloat}[!tb]
\begin{appfig}[diagrams (python)-type / clustered infrastructure, three clusters and one nested]
% Site cluster on the left, with the robot control network nested in its lowest
% row.  NAIRR resources upper right, public record lower right.  Accent tiles run
% on premises; gray tiles are shared or public.  Tile pitch 2.55 cm.
\dgnodew{dgtiled}{0.00}{0.00}{\glyphdb}{Data capture; the model cannot write}{m5}
\dgnodew{dgtiled}{2.55}{0.00}{\glyphai}{Advisory language model, on premises}{m1}
\dgnodew{dgtilem}{5.10}{0.00}{\glyphai}{Two verifier models, other developers}{m2}
\dgnode{dgtile}{0.00}{-2.25}{\glyphlock}{Segmentation: no route to control}{lk}
\dgnode{dgtile}{2.55}{-2.25}{\glyphdoc}{Audit trail, replayable}{m3}
\dgnode{dgtile}{5.10}{-2.25}{\glyphmon}{Surgeon's console, advice as text}{m4}
\dgnodew{dgtilem}{0.00}{-4.55}{\glyphrobot}{Eight-arm platform}{r1}
\dgnodew{dgtilem}{2.55}{-4.55}{\glyphstop}{Stops at 3 ms and 500 ms}{r2}
\dgnodew{dgtilem}{5.10}{-4.55}{\glyphgear}{Force limits, 3 N and 18 N}{r3}
\dgnodeg{dgtileg}{9.90}{0.00}{\glyphcloud}{Allocated compute, if granted}{n1}
\dgnodeg{dgtileg}{12.45}{0.00}{\glyphflask}{Synthetic cases from published values}{n2}
\dgnodeg{dgtileg}{9.90}{-2.25}{\glyphchart}{Measured disagreement rates}{n3}
\dgnodeg{dgtileg}{9.90}{-4.55}{\glyphlink}{Repository and identifiers}{p1}
\dgnodeg{dgtileg}{12.45}{-4.55}{\glyphdoc}{Report of rates and limits}{p2}
\begin{scope}[on background layer]
\node[inner sep=7pt,fit={(r1)(r1l)(r3)(r3l)}] (crx) {};
\node[dgcluster,fit={(m5)(m5l)(m2)(m2l)(crx)},inner sep=9pt] (cs) {};
\node[dgcluster2,fit={(r1)(r1l)(r3)(r3l)}] (cr) {};
\node[dgcluster2,fit={(n1)(n1l)(n2)(n2l)(n3)(n3l)}] (cn) {};
\node[dgcluster2,fit={(p1)(p1l)(p2)(p2l)}] (cp) {};
\end{scope}
\node[dgctitle,anchor=south west] at (cs.north west) {Clinical site, on premises: no patient data leaves it};
\node[dgctitle2,anchor=south east] at ([yshift=1pt]cr.north east) {Isolated robot control network};
\node[dgctitle2,anchor=south west] at (cn.north west) {NAIRR Pilot shared resources};
\node[dgctitle2,anchor=south west] at (cp.north west) {Public record};
\draw[dgedged] (m5) -- node[above,font=\tiny\sffamily,fill=fundwhite,inner sep=1pt] {read} (m1);
\draw[dgedgeb] (m1) -- node[above,font=\tiny\sffamily,fill=fundwhite,inner sep=1pt] {output} (m2);
\draw[dgedgeb] (m2l.south) -- (m4);
\draw[dgedge] (m2l.south west) -- (m3);
\draw[dgedged,-] (m1l.south west) -- (lk);
\draw[dgedged,-] (lkl.south) -- (r1);
\draw[dgedged] (m2) -- node[above,font=\tiny\sffamily,fill=fundwhite,inner sep=1pt] {code only} (n1);
\draw[dgedge] (n2) -- (n1);
\draw[dgedge] (n1l.south) -- (n3);
\draw[dgedge] (n3.east) -| (p2);
\draw[dgedge] (p2) -- (p1);
\legkey{-0.60}{-6.55}{fundaccent}{Runs on premises at the site}
\legkey{6.40}{-6.55}{fundgrayl}{Shared or public; synthetic data only}
\end{appfig}
\vspace{-0.60cm}
\figcaption{Figure 12. Where compute, data and review would sit, with the model and its verifiers\\
on premises, the robot network sealed off, and only code and synthetic cases shared.}
\end{appfloat}

\subsection{The quantitative record behind the robotic Phase 1}

A funder of the verification work is also, one step removed, a funder of the
trial it would make possible, and that trial rests on a dated quantitative
record. In August 2025 the company's ten-arm simulation of the agent returned a
median overall survival of 12.8 months against 5.4 for chemotherapy
\cite{qspmetpancre}. In May 2026 the developer's Phase 3 trial reported 13.2
months against 6.6 in the RAS G12 population of previously treated metastatic
disease \cite{rasolute302}, and the approval that followed reported 13.2 against
6.7 months in the overall population \cite{fdarasonque2026}. The company
describes the closeness of the simulated and reported figures as a chronology,
not a validation, because the populations, regimens and sample sizes differ.
Pancreatic ductal adenocarcinoma remains the cancer with the widest gap between
incidence and survival in the United States \cite{siegel2025statistics}, and
after resection the standard remains combination chemotherapy
\cite{conroy2018folfirinox}.

The program's direct cost is \$700,000 a year for five years, of which \$521,000
is personnel across 3.95 full-time equivalents \cite{movein2026}. The company's
comparable virtual trial work is \textbf{projected} at \$36,330 per run and about
one month, against industry benchmarks above \$120,000 and 4.5 months
\cite{capitalplan2026}. That figure describes simulation work, not the cost of
the clinical trial.

% sec-06: back matter.  Positioning, method, and the reference list under
% unsrturl, so every doi and url in references.bib reaches the printed page as a
% clickable target.
\section{Positioning, Method, and References}

\subsection{Positioning constraints}

An inquiry to apply is not an invitation to submit a proposal, an award, or an
assessment of eligibility, and nothing in this packet describes it as one. The
company holds no NSF award, and it is not eligible today for the
Commercialization Readiness Pilot or the Strategic Breakthrough tier
\cite{nsfncti2026,nsf26510}. No agreement of any kind exists with the National
Science Foundation, the University of Central Florida, the NAIRR Pilot, or the
National Institutes of Health, and none is described as a partner, sponsor, or
endorser. The Pitch does not ask NSF to fund a clinical trial, and no patient
data is used in the work it proposes.

Rasonque is approved for metastatic pancreatic adenocarcinoma after prior
systemic therapy \cite{fdarasonque2026}, and the perioperative use this program
proposes remains investigational. No investigational new drug application has
been submitted, and no patient has been treated. The private capital figure is
a plan, not a completed raise, and nothing in the day's capital instruction is
investment advice.

\subsection{Method and reproducibility}

This packet was produced from repository sources and from the public pages of
the offices it writes to: the tiers, amounts and gates from the current NSF
solicitations and the pilot's announcement \cite{nsf26510,nsf26511,nsfncti2026};
the split and the money frame from the capitalization plan
\cite{capitalplan2026}; the technical objectives from the deposited protocol and
investigational new drug application \cite{onpremwhippl,phase1ind}; the
verification practice from the company's three-model method
\cite{tripartisan2026}; and the staffing figures from the site package
\cite{movein2026}. Every recipient address was checked against the page its
office publishes, and every Pitch field was measured against its character limit
before it was committed.

The three figures are drawn in TikZ from the specifications in the day's
\appfile{diagrams} directory, one each in the D2, PlantUML and Diagrams
vocabularies, and no raster image is used anywhere. The human author retains
responsibility for every conclusion, and this document was adapted using Claude
Code Opus 5.5 \cite{repo2026}.

% Reserve room for the heading and the first lines of both columns together,
% so the heading is never stranded at the foot of a page.
\needspace{8\baselineskip}
\subsection{References}

\begin{apprefs}
\nocite{*}
\bibliographystyle{unsrturl}
\bibliography{references}
\end{apprefs}


\end{document}
